\honest{``Science'' as Curiosity-Stopper}{Eliezer Yudkowsky}{2007}

Suppose I conjured a brilliant light out of empty space on live television. Most people, I think, would be ``fairly curious.'' Now suppose I want to use this power in public without anyone becoming curious. I have a spell for that: whenever someone asks how I did it, I say ``Science!''

That word tells you nothing. It does not say whether the light will brighten or fade, or how to make one. The same trick works on an ordinary light bulb. In school the password is ``Electricity!'', and it does not let you calculate what to expect. \nb{``Curiosity-stopper'' and ``password'' are terms from the author's posts of the previous two weeks, ``Mysterious Answers to Mysterious Questions'' and ``Guessing the Teacher's Password.''}

If you thought a light bulb was scientifically inexplicable, it would take all of your attention. Being told that it is scientifically explicable means only, I say, ``that someone else knows how the light bulb works,'' and ``you don't know more than you knew earlier.'' \nb{The listener does learn something: that the bulb breaks no known law of nature and that the answer can be found. A bulb science could not explain would show that something the listener took as settled is mistaken. That is a reason to pay it more attention, whoever else knows.}

Someone will say that an unexplained bulb could bring fame and fortune. I am not talking about greed, I reply, but about the feeling of curiosity. Why should yours shrink because someone else knows? ``Is this not spite?'' \nb{The post gives no evidence of spite, and does not consider the reason in the previous note.}

If ``someone else knows'' could stop curiosity, then a voice saying ``Somebody already knows why this elephant is here'' would let you walk past a hovering green elephant. \nb{The elephant would still grip you because it fits nothing you know. That is the reason given two notes earlier.} I do not know where my kidneys are. I bet you do not know how your cerebellum, ``among other brain areas,'' coordinates a wave of your hand, so why should a light spell intrigue you more than raising your hand? I call this a ``double standard.''

The world is full of puzzles, I conclude. ``Prioritize, if you must.'' But do not complain that science has emptied the world of mystery. \nb{The reply is that much remains unknown to you. It says nothing about what you do come to understand. The author answered that part in ``Joy in the Merely Real,'' in 2008.}
